\subsection{Dataset and Evaluation Protocol}
\label{sec:dataset_and_metrics}

\paragraph{Dataset and Partitioning.}
We evaluated the detector on the KITTI 3D Object Detection dataset~\cite{geiger2012kitti}. The dataset contains $7{,}481$ calibrated LiDAR scans with synchronized rectified camera images and 3D bounding box annotations from road scenes. We used the 50/50 split introduced by Chen et al.~\cite{chen20153d}, with $3{,}712$ training frames and $3{,}769$ validation frames. Ground-truth sampling databases were generated exclusively from the $3{,}712$ training frames to prevent validation data leakage.

\paragraph{Spatial Region of Interest (ROI) and Target Parameterization.}
Following Section~\ref{sec:problem_formulation}, we aligned the LiDAR coordinate frame with the vehicle body (forward $X$, leftward $Y$, upward $Z$). We restricted perception to the Cartesian Region of Interest:
\begin{equation}
    \begin{aligned}
        x &\in [0.0, 70.4]\,\text{m}, \\
        y &\in [-40.0, 40.0]\,\text{m}, \\
        z &\in [-2.5, 1.0]\,\text{m}.
    \end{aligned}
\end{equation}
We applied an open-boundary inward margin of $\epsilon = 0.001\,\text{m}$ along all three axes. A horizontal grid resolution of $\Delta_x = \Delta_y = 0.1\,\text{m}$ produced an input tensor of $H_{\mathrm{in}} \times W_{\mathrm{in}} = 800 \times 704$ cells. An output stride of $s_{\mathrm{out}} = 4$ gave feature maps of size $H_{\mathrm{out}} \times W_{\mathrm{out}} = 200 \times 176$. We evaluated Car, Pedestrian, and Cyclist detections.

\paragraph{Evaluation Metrics.}
We evaluated oriented 2D Bird's-Eye View (BEV) footprints on the KITTI validation set using 40-point interpolated Average Precision (AP $R_{40}$). This metric averages interpolated precision at 40 recall positions $\mathcal{R}_{40} = \{1/40, 2/40, \dots, 1.0\}$:
\begin{equation}
    \mathrm{AP}|_{R_{40}} = \frac{1}{40} \sum_{r \in \mathcal{R}_{40}} \max_{r' \ge r} p(r').
\end{equation}
We used nominal oriented BEV Intersection-over-Union (IoU) thresholds of $\tau_{\mathrm{IoU}} = 0.70$ for Car and $\tau_{\mathrm{IoU}} = 0.50$ for Pedestrian and Cyclist. We grouped detections into three difficulty regimes:
\begin{itemize}
    \item \textbf{Easy}: 2D bounding box height $\ge 40\,\text{px}$, maximum occlusion level = 0 (fully visible), maximum truncation $\le 15\%$.
    \item \textbf{Moderate}: 2D bounding box height $\ge 25\,\text{px}$, maximum occlusion level = 1 (partially occluded), maximum truncation $\le 30\%$.
    \item \textbf{Hard}: 2D bounding box height $\ge 25\,\text{px}$, maximum occlusion level = 2 (difficult to see), maximum truncation $\le 50\%$.
\end{itemize}
Following autonomous driving literature, \textbf{Moderate macro mAP} ($\mathrm{mAP}_{\mathrm{M}}$) was adopted as the primary accuracy summary metric:
\begin{equation}
    \mathrm{mAP}_{\mathrm{M}} = \frac{1}{3} \sum_{c \in \{\mathrm{Car}, \mathrm{Ped}, \mathrm{Cyc}\}} \mathrm{AP}_{c, \mathrm{M}}.
\end{equation}
We also reported the unweighted mean across all nine class-difficulty combinations:
\begin{equation}
    \mathrm{mAP}_{9} = \frac{1}{9} \sum_{c \in \mathcal{C}} \sum_{d \in \{\mathrm{E}, \mathrm{M}, \mathrm{H}\}} \mathrm{AP}_{c, d}.
\end{equation}
E, M, and H denote Easy, Moderate, and Hard. $\mathrm{mAP}_{9}$ averages the nine class--difficulty AP values; it does not denote AP at an IoU threshold of 0.9.

\paragraph{Unified Benchmark Protocol and Evaluation Environment.}
The archived results use two evaluation implementations. MobilePIXOR and MobilePIXORNeXt use a local loader-aligned BEV evaluator with strict horizontal ROI filtering. PIXOR and the five external 3D detectors use a pinned reference evaluator with native ground-truth handling and preprocessing. Both report 40-point interpolation at the nominal class IoU thresholds, but this does not establish evaluator equivalence. Table~\ref{tab:comparison} therefore reports protocol groups without a shared AP ranking. Accuracy for the seed study was evaluated on an NVIDIA L4; the separate pipeline profiling snapshot for MobilePIXORNeXt and the external detectors used Tesla T4. Native ROI, checkpoint selection, training budgets, and model-stage boundaries also differ.
\subsection{Implementation Details}
\label{sec:implementation_details}

\begin{table}[t]
\centering\small
\caption{Experimental settings, architecture hyperparameters, and training configuration for MobilePIXORNeXt. Training uses a deterministic single-seed setup ($\text{seed} = 42$) with PyTorch 2.11.0 and CUDA 13.0. Timing uses a dedicated NVIDIA Tesla T4 GPU (16\,GB GDDR6 VRAM), batch size 1, and FP32 precision.}
\label{tab:setup}
\begin{tabular}{@{}>{\raggedright\arraybackslash}p{.38\columnwidth}>{\raggedright\arraybackslash}p{.56\columnwidth}@{}}
\toprule
\textbf{Configuration Element} & \textbf{Specification / Value} \\
\midrule
LiDAR Input Representation & Strata-BEV14 (14 channels; NumPy/Numba CPU) \\
Spatial ROI bounds ($X, Y, Z$) & $[0.00, 70.40]\,\text{m} \times [-40.00, 40.00]\,\text{m} \times [-2.50, 1.00]\,\text{m}$ \\
Grid resolution ($\Delta_x, \Delta_y, \Delta_z$) & $0.10\,\text{m} \times 0.10\,\text{m} \times 0.10\,\text{m}$ \\
Input / Output Grid Dimensions & $800 \times 704$ / $200 \times 176$ ($s_{\mathrm{out}} = 4$) \\
Detector Parameters & 753,569 trainable ($< 0.76\,\text{M}$); 9.34\,MB FP32 \\
Backbone Stages ($C_2, C_3, C_4$) & Widths: (48, 96, 128); Depths: (3, 4, 2) \\
Inverted Bottleneck Expansion & $e = 2.50$; $7 \times 7$ Depthwise Convolution \\
Attention Mechanism & Parameter-free Block-SimAM (9 residual blocks) \\
Multiscale Neck & Scale-Gated FPN (SG-FPN, $F_{32}$ fusion width) \\
Prediction Towers & GroupedHeader: Car ($K_1=1$), Ped/Cyc ($K_2=2$) \\
Loss Formulation & OGA (MGIoU + $\pi$-Corner) + IQA + T-SBUW \\
Data Augmentation & Physics-Consistent GT (PC-GTA) + Global \\
Optimizer / Base Learning Rate & AdamW (fused) / $\eta = 7.00 \times 10^{-4}$ \\
Minimum LR / Standard Schedule & $1.00 \times 10^{-6}$ / 50 epochs: 4 warmup + 46 cosine \\
Final Reported Checkpoint & Extended run: 160 epochs, 10 warmup; selected at epoch 130 \\
Weight Decay & $\lambda = 1.00 \times 10^{-3}$ (2D+ tensors only; biases/norm: 0.00) \\
Gradient Clipping Norm & Maximum $L_2$ norm = 10.00 \\
Batch Size (Train / Val) & 16 / 32 \\
Numerical Precision & Automatic Mixed Precision (Bfloat16 / FP32) \\
Hardware Environment & NVIDIA Tesla T4 GPU (16\,GB GDDR6, 98.44\,MiB peak) \\
\bottomrule
\end{tabular}
\end{table}

Table~\ref{tab:setup} summarizes the primary implementation settings and architectural hyperparameters across all pipeline components.

\paragraph{Network and Model Configuration.}
MobilePIXORNeXt processes 14-channel Strata-BEV14 tensors. We used a two-layer $3 \times 3$ convolutional stem ($32$ channels, stride 2) and three backbone stages with channel widths $(48, 96, 128)$ and block depths $(3, 4, 2)$. The stages used \texttt{MobilePIXORNeXt-S} blocks with $7 \times 7$ depthwise convolutions, an inverted expansion ratio of $e=2.5$, SiLU activations, and LayerScale initialized to $10^{-5}$. Parameter-free Block-SimAM attention was embedded in all nine residual blocks. The Scale-Gated Feature Pyramid Network (SG-FPN) fused stages C4, C3, and C2 into a $F_{32}$ (32-channel) representation at stride 4 ($200 \times 176$). Decoupled GroupedHeader towers processed Car ($K_1=1$) and Pedestrian/Cyclist ($K_2=2$) independently, with an initial $3 \times 3$ Conv-BN-SiLU transformation per group. The detector contained $753{,}569$ deployable parameters ($753{,}581$ trainable parameters during training including 12 transient T-SBUW multi-task logits), with an FP32 checkpoint size of $9.34\,\text{MB}$.

\paragraph{Training Schedule and Optimization.}
The configured optimizer is AdamW with base learning rate $\eta_0 = 7 \times 10^{-4}$, minimum learning rate $\eta_{\min} = 1.0 \times 10^{-6}$, weight decay $\lambda = 0.001$ on multidimensional weights, and maximum gradient $L_2$ norm $10.0$. Biases, normalization parameters, and loss-balancing logits are exempt from weight decay. Archived run metadata take precedence over a nominal common schedule: the final seed-42 run completed 160 epochs and its AP-selected checkpoint is at epoch 130, whereas final seeds 43 and 44 completed 100 epochs with selected checkpoints at epoch 100. The seed-42 MobilePIXOR run was archived at 38 of 100 planned epochs and evaluated at epoch 30; seeds 43 and 44 completed 100 epochs. The early B1/Binary35 run completed 100 epochs and used a loss-selected epoch-13 checkpoint. Most later historical component runs used a 50-epoch budget with AP-selected checkpoints, often before epoch 50. These records do not constitute a uniform-budget component or random-seed study. T-SBUW balances task losses as defined in Section~\ref{sec:losses_and_oga}.

\paragraph{Data Augmentation Pipeline.}
Training used Physics-Consistent Ground-Truth Augmentation (\texttt{PC-GTA}; Section~\ref{sec:pc_gta}). The final reported run enabled local ground fitting, static obstacle clearance, range-adaptive point subsampling, and visibility protection. The static obstacle clearance threshold was set to $N_{\mathrm{thresh}} = 2$ points within the expanded footprint, and estimated ground elevation was bounded within $[-2.5, -0.8]\,\text{m}$. LOS rejection, shadow masking, and radiometric calibration were available as configurable options but disabled in the primary benchmark. Ground height was estimated within $R(d) = \operatorname{clip}(2.5 + 0.03\max(0, d - 15), 2.5, 4.0)\,\text{m}$. GT sampling was applied with probability $0.5$. Subsequent world transforms included flipping along the $X$-axis with probability $0.5$, rotation $\theta \sim \mathcal{U}(-45^\circ,45^\circ)$, and isotropic scaling $s \sim \mathcal{U}(0.95,1.05)$. World translation was applied with probability $0.5$, using independent Gaussian noise with standard deviations $(0.2,0.2,0.1)\,\text{m}$ along $(X,Y,Z)$.

\begin{table*}[t]
\centering\small
\caption{Quantitative object detection comparison on the KITTI validation set for Bird's-Eye View (BEV) bounding box estimation. Average Precision (AP, \%) is evaluated using 40-point interpolation ($R_{40}$) at Moderate difficulty ($\text{IoU} \ge 0.70$ for Car; $\text{IoU} \ge 0.50$ for Pedestrian and Cyclist). The top partition compares compact and real-time BEV/pillar detectors; the bottom partition lists heavier two-stage 3D anchor-based detectors as external reference points. Best values among compact models are in \textbf{bold}, and second-best values are \underline{underlined}.}
\label{tab:comparison}
\setlength{\tabcolsep}{3pt}
\begin{tabular}{@{}lrrrrrr@{}}
\toprule
\textbf{Method} & \textbf{\shortstack{Params\\(M) $\downarrow$}} & \textbf{\shortstack{Car AP$_{\mathrm{M}}$\\(\%) $\uparrow$}} & \textbf{\shortstack{Pedestrian\\AP$_{\mathrm{M}}$ (\%) $\uparrow$}} & \textbf{\shortstack{Cyclist\\AP$_{\mathrm{M}}$ (\%) $\uparrow$}} & \textbf{\shortstack{mAP$_{\mathrm{M}}$\\(\%) $\uparrow$}} & \textbf{\shortstack{mAP$_{9}$\\(\%) $\uparrow$}} \\
\midrule
PIXOR (Car only)~\cite{yang2018pixor} & 2.135 & 72.02 & -- & -- & -- & -- \\
MobilePIXOR~\cite{nguyen2024indoor} & \textbf{0.598} & 75.52 & 28.16 & 20.53 & 41.40 & 42.27 \\
PointPillars~\cite{lang2019pointpillars} & 4.835 & \underline{88.02} & 56.01 & 66.34 & 70.13 & 72.25 \\
SECOND~\cite{yan2018second} & 5.326 & \textbf{88.56} & \underline{56.56} & \textbf{71.16} & \textbf{72.09} & \textbf{73.79} \\
\textbf{MobilePIXORNeXt (Ours)} & \underline{0.754} & 86.90 & \textbf{60.75} & \underline{66.73} & \underline{71.46} & \underline{73.42} \\
\midrule
\multicolumn{7}{@{}l}{\textit{Heavyweight 3D Anchor-Based Reference Detectors}} \\
PointRCNN~\cite{shi2019pointrcnn} & 4.039 & 88.42 & 57.84 & 73.87 & 73.38 & 75.31 \\
Part-A$^2$-Anchor~\cite{shi2020part} & 63.815 & 90.05 & 64.20 & 74.60 & 76.28 & 78.04 \\
PV-RCNN~\cite{shi2020pvrcnn} & 13.116 & 90.29 & 58.48 & 74.54 & 74.44 & 76.49 \\
\bottomrule
\end{tabular}
\par\vspace{2pt}
{\footnotesize \textit{Note.} $\text{mAP}_{\mathrm{M}}$ is the unweighted mean across the three target classes at Moderate difficulty. $\text{mAP}_{9}$ denotes the macro-average over all nine class-difficulty combinations (Easy, Moderate, Hard). Dashes (--) indicate classes unsupported by single-class PIXOR. All results use the same KITTI validation split. MobilePIXORNeXt and historical MobilePIXOR use local loader-aligned BEV AP; PIXOR, PointPillars, SECOND, and the heavier references report published reference KITTI AP. Under our unified local evaluation harness (Figure~\ref{fig:pr_curves}), PointPillars achieves $55.83\%$ and SECOND achieves $56.47\%$ on Pedestrian Moderate AP. GPU timings are specified separately in Table~\ref{tab:efficiency}.}
\end{table*}


\paragraph{Batching, Precision, and Hardware Environment.}
The configured training batch size was 16, with validation batch size 32 and six persistent DataLoader workers. The saved seed-study reports record FP32 inference on NVIDIA L4 with PyTorch 2.11.0 and CUDA 13.0. Offline timing uses batch size 1. The separate T4 profiling report records $98.44\,\text{MiB}$ peak allocated evaluation memory. Architecture, hardware, checkpoint epoch, and evaluation protocol must be matched before using archived timings or AP differences to attribute a component gain. Table~\ref{tab:seed_robustness} summarizes the available runs and explicitly identifies their configuration and budget differences.
